\documentclass[10pt,a4paper]{amsart}
\usepackage[margin=19mm]{geometry}
\usepackage{amsmath,amssymb,mathtools}
\usepackage[scr=rsfs,cal=euler]{mathalfa}
\usepackage{xfrac}
\usepackage[unicode,hidelinks,bookmarks=false]{hyperref}
\newcommand{\ie}{i.e.\ }
\newcommand{\cf}{cf.\ }
\newcommand{\resp}{respectively\ }
\newcommand{\Subsection}{\subsection}
\providecommand{\nobreakdash}{\nobreak\hbox{-}}
\setlength{\emergencystretch}{2em}
\multlinegap0pt
\usepackage{amssymb}
\usepackage{bm}
\usepackage{mathtools}
\usepackage{xcolor}
\usepackage{enumerate}
\newtheorem{lem}{Lemma}
\newtheorem{thm}{Theorem}
\newcommand{\AAA}{\mathrm {A}}
\newcommand{\Alt}{\mathrm {Alt}}
\newcommand{\CB}{\mathbf{C}}
\newcommand{\Irr}{\mathrm {Irr}}
\newcommand{\OB}{\mathbf{O}}
\newcommand{\PSp}{\mathrm {S}}
\newcommand{\SSS}{\mathrm {S}}
\newcommand{\Sym}{\mathrm {Sym}}
\newcommand{\ZB}{\mathbf{Z}}
\renewcommand{\geq}{\geqs}
\newcommand{\geqs}{\geqslant}
\newcommand{\la}{\langle}
\renewcommand{\leq}{\leqs}
\newcommand{\leqs}{\leqslant}
\newcommand{\ra}{\rangle}
\newcommand{\supp}{\mathrm {supp}}
\title{Orders of commutators and Products of conjugacy classes in finite groups}
\author{Hung P. Tong-Viet}
\date{Source: arXiv:2507.10882v3 (CC BY 4.0)}
\begin{document}
\maketitle

\noindent\emph{Source and licence.} Orders of commutators and Products of conjugacy classes in finite groups. Version arXiv:2507.10882v3 (CC BY 4.0), distributed by arXiv under the Creative Commons Attribution 4.0 International licence, http://creativecommons.org/licenses/by/4.0/, as recorded in the arXiv licence metadata for this exact version and shown on the abstract page https://arxiv.org/abs/2507.10882v3. Full licence notice: \texttt{licenses/arxiv-2507.10882-CC-BY-4.0.txt}.\par\smallskip

\noindent\textit{Editorial scope.} Counted unit: the article's thm th:almost-simple (source0.tex lines 306--308). The article's complete original proof of it is delivered with it (source0.tex lines 310--357, one \texttt{proof} environment of 48 lines). No mathematics is written, repaired or re-derived: the statement and the proof are the article's own lines, in the article's own notation, and no retained line is substituted or re-flowed.  Retained uncounted as the source-local context the proof needs: lines 228--230; lines 250--252; lines 282--284.  Nothing was omitted from any retained span, so the package carries no omission record.  The retained lines cite 4 works, whose entries are transcribed verbatim from the article's own bibliography source in the e-print and delivered with short editorial labels recorded in the item's reference mapping.  No retained line contains a figure environment, an image-inclusion command or drawing code, so no image is delivered and the package declares no asset dependency.  Editorial formatting only, in addition: an AMS \texttt{amsart} preamble that restates the article's own declarations and macros used by the retained lines, namely the environments \texttt{lem}, \texttt{thm} on separate counters, together with the standard packages those lines need.  The retained environments print in this document as Lemma 1, Theorem 1 in order of appearance, whereas the article prints the same items with the numbering of its own full document.  The complete native source of the article as distributed in the arXiv e-print for this exact version is bundled unchanged as \texttt{sources/181589/source0.tex}.
\begin{lem}\label{lem:pa}
Let $G$ be a finite group and let $x\in G.$ If $|x^G|$ is a power of a prime, then $\la x^G\ra$ is a solvable group.
\end{lem}

\begin{lem}\label{lem:semisimple}
Let $S$ be a finite nonabelian simple group of Lie type in defining characteristic $r$. For each prime $p\neq r$ dividing $|S|$, $S$ has an irreducible character of $p$-defect zero whose degree is not divisible by $r$.
\end{lem}

 \begin{lem}\label{lem:Sylow-centralizers}
 Let $G$ be a finite almost simple group with socle $S$, where $S$ is  a finite simple group of Lie type in characteristic $r$.  Let $R$ be a Sylow $r$-subgroup of $G$. Then $\CB_G(R)\leq R.$ 
 \end{lem}

\begin{thm}\label{th:almost-simple}
Let $G$ be a finite almost simple group with simple socle $S$. Let $p$ be a prime and let $x\in G.$ If $[x,g]=1$ or $[x,g]$ is $p$-singular for every $g\in G$, then $x=1$.
\end{thm}

\begin{proof}
Let the pair $(G,x)$ be a counterexample to the theorem with $|G|$ minimal. Then $1\neq x\in G$ and $[x,g]=1$ or  $[x,g]$ is $p$-singular for every $g\in G.$ Suppose that $H$ is a finite almost simple proper subgroup of $G$ containing $x$. Clearly, $[x,h]=1$ or it is $p$-singular for every $h\in H.$  By the minimality of $|G|$, $x=1$, which is a contradiction. Thus $x$ does not lie in any almost simple proper subgroup of $G$. In particular,  $G=\la S,x\ra=S\la x\ra$ and hence $G'=S$. Let $K=x^G.$ Then 
\begin{equation}\label{eqn1} K^{-1}K=C_0\cup C_1\cup C_2\cup\dots\cup C_a
\end{equation} where $C_0=1$ and for $1\leq i\leq a,$  $C_i$ is a conjugacy class of $G$ containing $y_i$ with $y_i\in G$ a nontrivial $p$-singular element. Taking the class sums, we get 
\begin{equation}\label{eqn2}\widehat{K^{-1}}\widehat{K}=\sum_{i=0}^a n_i \widehat{C_i}\end{equation}
 where the structure constant $n_i=n(K^{-1},K,C_i)$ can be computed by \begin{equation}\label{eqn3}n(K^{-1},K,C_i)=\dfrac{|K|^2}{|G|}\sum_{\phi\in\Irr(G)}\dfrac{|\phi(x)|^2\overline{\phi(y_i)}}{\phi(1)}\end{equation} (see Exercise (3.9) in \cite{Isaacs-ct}). For $i=0,$ we have $n_0=|K|.$
Let $\chi\in\Irr(G)$. Applying the homomorphism $\omega_\chi$ to the class sum Equation \eqref{eqn2}, we get

\begin{equation}\label{eqn4}
\dfrac{|K|^2|\chi(x)|^2}{\chi(1)^2}=|K|+\sum_{i=1}^a n_i\dfrac{|C_i|\chi(y_i)}{\chi(1)}.
\end{equation}

(a) Assume first that $S$ is a sporadic simple group, the Tits group or an alternating group $\AAA_n$ with $5\leq n\leq 19$. For these cases, the character tables of the almost simple groups $G$ with socle $S$ are available in GAP \cite{GAP}. For each $1\neq x\in G,$ we can compute the structure constant $n(K^{-1},K,C)$ for the conjugacy classes $C=y^G$ for every $y\in G$, where $K=x^G,$ and one can check that there exist two nontrivial elements $z_1,z_2$ of coprime orders with nonzero structure constants. Hence $z_1^G$ and $z_2^G$ appear in the decomposition of the product $K^{-1}K$, which is a contradiction.

(b) Assume next that $S\cong \AAA_n$ with $n\ge 20.$ For these cases, $G\in \{\AAA_n,\SSS_n\}.$ Let $\Omega=\{1,2,\ldots, n\}$. If $\sigma=(i_1,i_2,\dots,i_s)$ is an $s$-cycle, then we denote by $\ell(\sigma)$ the length of $\sigma$, which is $s$, and the support of $\sigma$ is the subset $\supp(\sigma)=\{i_1,i_2,\ldots,i_s\}$ of $\Omega.$ Let $x=\sigma_1\sigma_2\dots\sigma_k$, where $\sigma_1, \dots,\sigma_k$ are pairwise disjoint  cycles in $\SSS_n$, including cycles of length $1$. Denote $\ell_i=\ell(\sigma_i)$ for the length of each cycle, and assume without loss of generality that  $\ell_1\geq \ell_2\geq\dots\geq \ell_k.$ Note that $\sum_{i=1}^k \ell_i=n.$

Assume $x$ has a fixed point on $\Omega$. Then $\ell_k=1$. Without loss of generality, assume $\sigma_k=(n)$. Then $x=\sigma_1\sigma_2\dots\sigma_{k-1}\in G\cap \Sym(\Omega_1)$ with $\Omega_1=\Omega\setminus \{n\}$. Now $G\cap \Sym(\Omega_1)$ is a finite almost simple group with socle $\Alt(\Omega_1)\cong \AAA_{n-1}$ containing $x$, which is impossible.  Thus  $x$ has no fixed point and hence $\ell_k\ge 2.$ 




Assume  $G=\SSS_n$. Suppose that $k\ge 2.$ Then $\ell_k\leq n/2$. Let $x=uv$, where $v=\sigma_k$ and $u=\sigma_1\sigma_2\dots \sigma_{k-1}$. Then $1\neq u\in\Sym(\Omega_2)$ and $v\in\Sym(\Omega_1)$ with $\Omega_1=\supp(v)$ and $\Omega_2=\Omega\setminus \Omega_1.$  Here $n>m:=|\Omega_2|=n-\ell_k\ge n/2\ge 10$. Let $g\in\Sym(\Omega_2)\leq G$, we see that $[g,v]=1$ and thus $[x,g]=[uv,g]=[u,g]$ is either $1$ or $p$-singular. Thus $(\Sym(\Omega_2),u)$ satisfies the hypothesis of the theorem. By the minimality of $|G|$, we have $u=1$,  a contradiction. Suppose that $k=1$ so $x$ is an $n$-cycle. Without loss of generality, assume $x=(1,2,\dots,n)$. Let $g_1=(1,3,2)$ and $g_2=(1,2)(3,4)$. Then both $g_1$ and $g_2$ lie in $\AAA_n\unlhd G$. By computation, we get $[x,g_1]=(1,3,n)$ and $[x,g_2]=(1,3,4,2,n)$, violating the hypothesis of the theorem.

Assume $G=\AAA_n.$ Observe that if $k=1$, then $x\in G$ is an $n$-cycle and $n$ must be odd. The argument above also shows that this is impossible. Thus, we assume $k\ge 2.$  Then $2\leq \ell_k\leq n/2$. Let $\Omega_1$ be the support of $\sigma_k$ and $\Omega_2=\Omega\setminus\Omega_1$. Then $m=|\Omega_2|=n-\ell_k\ge  n/2 \ge 10.$ Let $H=\Alt(\Omega_2)\cong \AAA_m$. Then $H$ is a finite nonabelian simple subgroup of $G$ and $H$ centralizes $v=\sigma_k$. Note that $x=uv=vu\in G.$

Assume $\ell_k$ is odd. Then $v$ is an even permutation. So $v\in\Alt(\Omega_1)$ and $u\in \Alt(\Omega_2)\leq G$. Let $g\in  H$. Then $[g,v]=1$ and thus $[x,g]=[u,g]$ is either 1 or $p$-singular and thus $u=1$ as $|H|<|G|$, a contradiction.

 Assume $\ell_k$ is even. Then $v$ is an odd permutation and so $u\in\Sym(\Omega_2)$ is odd as $x=uv$ is an even permutation.  Let $g\in\Sym(\Omega_2)$. If $g$ is an even permutation, then $g\in\Alt(\Omega_2)\leq G$; and if $g$ is an odd permutation, then $gv\in G$. As $[g,v]=1$, if $g$ is even, then $[u,g]=[x,g]$; and if $g$ is odd, then $[x,gv]=[uv,gv]=[u,g]$. In either cases, $[u,g]$ is either $1$ or it is $p$-singular for all $g\in \Sym(\Omega_2)$. Hence the pair $(\Sym(\Omega_2),u)$ satisfies the hypothesis of the theorem with $|\Sym(\Omega_2)|<|G|$. Thus $u=1$, which is a contradiction.

(c) Assume that $S\not\in\{\PSp_4(2)', {}^2\textrm{F}_4(2)'\}$ is a finite simple group of Lie type in characteristic $r$, where $r$ is a prime.  Note that the cases $S\cong \PSp_4(2)'\cong \AAA_6$  and $S\cong {}^2\textrm{F}_4(2)'$ have been handled earlier. 

Assume first that $p=r$. Let $\theta\in\Irr(S)$ be the Steinberg character of $S$ of degree $|S|_p$. It is well-known that (see \cite{F1}) $\theta$ extends to $\chi\in\Irr(G)$ and that $\theta(y)=0$ whenever $y\in S$ is $p$-singular since $\theta$ has $p$-defect zero (\cite[Theorem 8.17]{Isaacs-ct}). For each $i$ with $1\leq i\leq a$, we know that $y_i=[x,g_i]$ for some $g_i\in G$. As $G/S$ is abelian,  $y_i\in G'=S$ for all $i\ge 1$. Moreover, $y_i$ is $p$-singular by the hypothesis. Therefore $\chi(y_i)=\theta(y_i)=0$ for all $1\leq i\leq a$. From Equation \eqref{eqn4}, we obtain 
\begin{equation}\label{eqn5}
|\chi(x)|^2=\frac{\chi(1)^2}{|K|}=\frac{|S|_p^2}{|K|}.
\end{equation}
By \cite[Corollary 3.6]{Isaacs-ct}, $\chi(x)$ is an algebraic integer and so is $|\chi(x)|^2=\chi(x)\chi(x^{-1})$. Now the previous equation implies that $|\chi(x)|^2$ is also a rational integer and so it is an integer by \cite[Lemma 3.2]{Isaacs-ct}. It follows that $|K|$ divides $|S|_p^2$. In particular, $|K|$ is a power of $p$. By Lemma \ref{lem:pa}, $\la K\ra=\la x^G\ra$ is solvable. However, since $G$ is almost simple, $G$ has no nontrivial solvable normal subgroup. Thus $\la x^G\ra=1$ forcing $x=1$, a contradiction.


Assume that $p\neq r.$ By Lemma \ref{lem:semisimple}, $S$ has an irreducible character $\theta$ of $p$-defect zero and $r\nmid \theta(1)$.  Let $\chi\in\Irr(G)$ be an irreducible character of $G$ lying above $\theta$. By Clifford's theorem (\cite[Theorem 6.2]{Isaacs-ct}), $\chi_S=e(\theta_1+\dots+\theta_t)$, where each $\theta_i$ is conjugate to $\theta=\theta_1$ and $t=|G:I_G(\theta)|$, here $I_G(\theta)$ is the inertia group of $\theta$ in $G$.  Since $\theta\in\Irr(S)$ has $p$-defect zero, each $\theta_i$ also has $p$-defect zero and thus $\theta_i(y)=0$ for every $p$-singular element $y\in S.$ Arguing as in the previous case, we know that $y_j\in S$ and it is $p$-singular  for all $j$ with $1\leq j\leq a$. Thus $\theta_i(y_j)=0$ for every $i$ and $j$ with $1\leq i\leq t$ and $1\leq j\leq a.$ It follows that $\chi(y_j)=e\sum_{i=1}^t\theta_i(y_j)=0$ for all $j$ with $1\leq j\leq a$. Equation \eqref{eqn4} yields
\begin{equation}\label{eqn6}
|\chi(x)|^2=\frac{\chi(1)^2}{|K|}.
\end{equation}
We claim that $\chi$ is an extension of $\theta$.  If $t>1$, then $\chi=\varphi^G$ for some $\varphi\in\Irr(I_G(\theta))$ lying over $\theta$. In this case, $S\unlhd I_G(\theta)\unlhd G$ and $x\not\in I_G(\theta)$ which implies that $\chi(x)=\varphi^G(x)=0$. However, this is impossible by Equation \eqref{eqn6}. Therefore $t=1$ and so $\theta$ is $G$-invariant. Since $G/S$ is cyclic, by \cite[Corollary 11.22]{Isaacs-ct} $\theta$ extends to $G$ and thus $\chi$ is an extension of $\theta$. Hence $\chi(1)=\theta(1)$ is not divisible by $r$. Recall that $r$ is the characteristic of $S$.

Arguing as in the previous case, $|\chi(x)|^2$ is both an algebraic integer and a rational integer, it is an integer and thus $|K|$ divides $\chi(1)^2$. Hence $r$ does not divide $|K|=|G:\CB_G(x)|$. Thus $\CB_G(x)$ contains a Sylow $r$-subgroup $R$ of $G$. In other words, $x\in \CB_G(R)$. By Lemma \ref{lem:Sylow-centralizers}, $x\in R$. In particular, $x$ is a nontrivial $r$-element.    
 
 Suppose that $x^G\cap \CB_G(x)=\{x\}$. Since $x$ is an $r$-element, by Glauberman $\ZB_p^*$-theorem (\cite[Theorem 5.1]{GTT}), $x\in \ZB_r^*(G)$. However, as $G$ is an almost simple group with simple socle $S$, a finite simple group of Lie type in characteristic $r$, $\OB_{r'}(G)=1=\ZB(G)$. Hence $x\in \ZB_r^*(G)=1$, a contradiction.  Thus we can find $g\in G$ such that $x^g\neq x$ and $[x,x^g]=1$. Hence $[x,g]=x^{-1}x^g$ is a nontrivial $r$-element. On the other hand, $[x,g]$ is $p$-singular. Thus $p=r$, contradicting our assumption that $p\neq r.$  The proof is now complete.
\end{proof}

\begin{thebibliography}{99}
\bibitem[F1]{F1} W. Feit, Extending Steinberg characters, in { Linear algebraic groups and their representations (Los Angeles, CA, 1992)}, 1--9, Contemp. Math., 153, Amer. Math. Soc., Providence, RI, 1993.
\bibitem[GAP]{GAP} The GAP Group, {GAP -- Groups, Algorithms, and Programming, Version 4.7.5}, 2014, \texttt{(http://www.gap-system.org)}.
\bibitem[GTT]{GTT} R.~M. Guralnick, H.~P. Tong-Viet\ and\ G.~M. Tracey, Weakly subnormal subgroups and variations of the Baer-Suzuki theorem, J. Lond. Math. Soc. (2) {\bf 111} (2025), no.~1, Paper No. e70057, 29 pp.
\bibitem[Isaacs]{Isaacs-ct} I.~M. Isaacs, {Character theory of finite groups}, AMS Chelsea Publ., Providence, RI, 2006.
\end{thebibliography}
\end{document}
